\begingroup%
\def\baselinestretch{1}%
\definecolor{D09Blue}{HTML}{3478B8}%
\definecolor{D09Teal}{HTML}{27877F}%
\definecolor{D09Amber}{HTML}{C58B28}%
\definecolor{D09Ink}{HTML}{243447}%
\definecolor{D09PaleBlue}{HTML}{F0F6FC}%
\definecolor{D09PaleAmber}{HTML}{FFF8E8}%
\begin{tikzpicture}[
  x=1mm, y=1mm,
  card/.style={draw, fill=white, rounded corners=2mm, line width=.75pt,
    inner xsep=2.5mm, inner ysep=2.6mm, outer sep=.4pt, align=flush center,
    text=D09Ink, font=\fontsize{9.5}{12.5}\selectfont},
  group box/.style={draw, rounded corners=2.5mm, line width=.7pt,
    inner sep=0pt, outer sep=.4pt},
  group title/.style={inner sep=0pt, text=D09Ink, align=flush center,
    font=\bfseries\fontsize{10}{12.5}\selectfont},
  edge text/.style={fill=white, inner xsep=1mm, inner ysep=.5mm, align=flush center,
    text=D09Ink, font=\fontsize{9}{12.5}\selectfont},
  flow/.style={draw=D09Ink, line width=.7pt, rounded corners=1.8mm,
    -{Latex[length=2.5mm,width=1.6mm]}},
  schema flow/.style={flow, draw=D09Blue, line width=.9pt,
    -{Latex[length=3mm,width=2mm]}}
]
\node[group title] at (0,17) {原生 AKA 模式：单层混杂循环（无高层约束与方向锚点）};
\node[group title] at (0,-48)
  {我们的两阶段解耦模式：Planner（管压缩态设计）+ Coder（管解压态恢复）};

\node[card, draw=D09Amber, minimum width=42mm, minimum height=18mm] (o1) at (-55,0)
  {从零起步 / 拉取已有 SOTA};
\node[card, draw=D09Amber, minimum width=46mm, minimum height=18mm] (o2) at (0,0)
  {单一 Episode 混杂迭代\\
   边猜策略 · 边改底层代码 ·\\
   边排编译错};
\node[card, draw=D09Amber, minimum width=46mm, minimum height=18mm] (o3) at (55,0)
  {验证与晋升判定\\
   易因局部实现摩擦\\
   放弃正确方向};

\node[card, draw=D09Blue, minimum width=51mm, minimum height=24mm] (n1) at (-48,-71)
  {Planner 维护 Schema\\
   Design / Extract 建立初版\\
   Optimize 分层级演进设计};
\node[card, draw=D09Teal, minimum width=55mm, minimum height=24mm] (n2) at (48,-71)
  {Coder 维护 Kernel（改造自 AKA）\\
   Implement 定向信息解压\\
   与代码恢复\\
   内部消化编译与工程细节};

\begin{scope}[on background layer]
  \node[group box, draw=D09Amber, fill=D09PaleAmber,
    fit={(-81,23) (81,-25)}] (old) {};
  \node[group box, draw=D09Blue, fill=D09PaleBlue,
    fit={(-81,-41) (81,-88)}] (new) {};
\end{scope}

\draw[flow] (o1) -- (o2);
\draw[flow] (o2) -- (o3);
\coordinate (retry) at (0,-18);
\draw[flow, draw=D09Amber] (o3.south) -- (o3.south |- retry)
  -- node[edge text, fill=D09PaleAmber] {无方向盲目重试} (retry) -- (o2.south);

\coordinate (schema) at (0,-65);
\coordinate (evidence) at (0,-77);
\draw[schema flow] (n1.east |- schema) --
  node[edge text, fill=D09PaleBlue] {下发固定 Schema\\明确方向与硬边界} (n2.west |- schema);
\draw[schema flow, draw=D09Teal] (n2.west |- evidence) --
  node[edge text, fill=D09PaleBlue] {回传干净实测证据\\不稀释 Planner 注意力} (n1.east |- evidence);

\draw[flow, draw=D09Blue, dash pattern=on 2.6pt off 2pt] (old.south) --
  node[edge text, right=2mm] {注入 Kernel Schema 拆分职责} (new.north);
\end{tikzpicture}%
\endgroup%
